% !PRINTED 246
\weberpage{246}
%<<CONT
in the ordinary sense. All numbers at which one arrives in this way can be brought to the form

\weq{2}{\alpha = a_0 + a_1\varepsilon + \cdots + a_{n-1}\varepsilon^{n-1}}

Since in calculating with such numbers the modulus $p$ is always held fixed, we shall designate two numbers that are congruent with respect to the modulus $p$ simply as \gesperrt{equal}. The objects of the calculation are then properly not the individual numbers themselves, but the \gesperrt{classes of numbers} consisting of all numbers congruent to one another. These numbers $\alpha$ are called the \gesperrt{Galois imaginaries}. The system of such imaginary numbers belonging to a definite $\varepsilon$ we denote for brevity by $\mathfrak{C}$.

We then have first the theorem:

\gesperrt{2. There are $p^n$ and no more distinct numbers in $\mathfrak{C}$.}

But from Theorem 1. it follows further:

\gesperrt{3. The product of two or more numbers from $\mathfrak{C}$ is equal to zero if and only if at least one of the factors is equal to zero.}

We obtain the complete system of numbers $\mathfrak{C}$ if the rational numbers $a_0, a_1, \ldots, a_{n-1}$ in (2) are each made to run through a complete system of residues with respect to the modulus $p$. Every system $\mathfrak{C}$ contains the residues of the natural numbers $0, 1, \ldots, p-1$, and for $n = 1$, $\mathfrak{C}$ is identical with this system.

If $\alpha$ is a fixed number in $\mathfrak{C}$ different from zero, then $\alpha\xi$ runs through the full system $\mathfrak{C}$ together with $\xi$. For it follows from 3. that $\alpha\xi$ can be equal to $\alpha\xi'$ only when $\xi = \xi'$. From this it follows:

\gesperrt{4. If $\alpha, \beta$ are two given numbers in $\mathfrak{C}$ and $\alpha$ is different from zero, then there is one and only one number $\gamma$ in $\mathfrak{C}$ for which the condition}

\[
\alpha\gamma = \beta
\]

\gesperrt{is satisfied.}

With this the operation of division in the system $\mathfrak{C}$ is also shown to be permissible. We denote the number $\gamma$ whose existence is asserted in 4. by

\[
\beta : \alpha \quad \text{or} \quad \frac{\beta}{\alpha}.
\]
